\documentclass[12pt,a4paper,oneside]{report}
% A generic thesis skeleton. Most institutions mandate their own class or
% style file — when yours does, drop it beside this file and swap the line
% above. The structure below (front matter, numbered chapters, appendix,
% bibliography) is what nearly every regulation asks for.

\usepackage[T1]{fontenc}
\usepackage[utf8]{inputenc}
\usepackage[margin=2.5cm]{geometry}
\usepackage{amsmath, amssymb}
\usepackage{graphicx}
\usepackage{booktabs}
\usepackage{setspace}
\usepackage[hidelinks]{hyperref}

\onehalfspacing

\title{Your Thesis Title}
\author{Your Name}
\date{\today}

\begin{document}

% ---------- Front matter ----------
\pagenumbering{roman}
\maketitle

\chapter*{Abstract}
\addcontentsline{toc}{chapter}{Abstract}
A concise summary of the whole thesis: problem, approach, findings,
significance. Usually one page.

\chapter*{Acknowledgements}
\addcontentsline{toc}{chapter}{Acknowledgements}
Thank the people and funders who made the work possible.

\tableofcontents
\listoffigures
\listoftables

% ---------- Main matter ----------
\cleardoublepage
\pagenumbering{arabic}

\chapter{Introduction}
\label{ch:introduction}

Motivate the research, state the questions, and outline the thesis
structure. Reference prior work as you go~\cite{knuth1984}.

\chapter{Literature Review}
\label{ch:literature}

Organise by theme, and end with the gap your work addresses.

\chapter{Methodology}
\label{ch:methodology}

\begin{equation}
  \label{eq:example}
  \mathcal{L}(\theta) = -\frac{1}{N}\sum_{i=1}^{N} \log p_\theta(y_i \mid x_i) .
\end{equation}

\chapter{Results}
\label{ch:results}

\begin{table}[ht]
  \centering
  \begin{tabular}{lrr}
    \toprule
    Condition & Measure A & Measure B \\
    \midrule
    Control   & 0.42      & 1.8       \\
    Treatment & 0.67      & 2.4       \\
    \bottomrule
  \end{tabular}
  \caption{Replace with your own results.}
  \label{tab:results}
\end{table}

\chapter{Discussion}
\label{ch:discussion}

Interpret the results against the questions from
Chapter~\ref{ch:introduction}, and be explicit about limitations.

\chapter{Conclusion}
\label{ch:conclusion}

Restate the contribution and outline future work.

% ---------- Back matter ----------
\appendix
\chapter{Supplementary Material}
\label{app:supplementary}

Extra derivations, full tables, instrument text.

\bibliographystyle{plain}
\bibliography{references}
\addcontentsline{toc}{chapter}{Bibliography}

\end{document}
